\chapter{Akathisia}

\section*{Definition}
Akathisia is a movement disorder characterized by distressing inner restlessness and an irresistible urge to move, often accompanied by pacing, rocking, shifting weight, or repeatedly moving the legs. It is most often caused by medicines that affect dopamine, especially antipsychotics, but can have other causes. It may begin soon after a medicine is started or increased, or persist after longer exposure; the timing helps guide evaluation but does not by itself establish the diagnosis.

\section*{When to See a Doctor}

\begin{itemize}
 \item Subjective inner restlessness developing after starting, increasing, or switching a medicine that can cause akathisia
 \item Objective fidgeting, pacing, or inability to sit still
 \item Distress severe enough to affect sleep, work, or social function
 \item Akathisia with suicidal thoughts, self-harm urges, or thoughts of harming someone else -- seek emergency help immediately
\end{itemize}

\section*{How It Develops}
The mechanism is not fully understood. Changes in dopamine signaling and interactions with serotonin and norepinephrine pathways may disturb circuits involved in movement, attention, and emotional regulation. Akathisia is more common with some high-potency or rapidly increased antipsychotics, but any antipsychotic can cause it. Antidepressants, antiemetics, and other medicines can also trigger it. Iron deficiency may contribute in some people, although it is not a complete explanation for medication-induced akathisia.

\section*{Causes}
\begin{itemize}
 \item \textbf{Antipsychotics:} First- and second-generation drugs, including haloperidol, risperidone, olanzapine, and aripiprazole; risk varies by drug, dose, and rate of increase
 \item \textbf{Antidepressants:} SSRIs, SNRIs, tricyclic antidepressants, MAO inhibitors
 \item \textbf{Antiemetics:} Metoclopramide, prochlorperazine, droperidol
 \item \textbf{Calcium channel blockers:} Flunarizine and cinnarizine, where these medicines are available
 \item \textbf{Dopamine-depleting agents:} Tetrabenazine, reserpine
 \item \textbf{Other:} Lithium, buspirone, and withdrawal from benzodiazepines or other sedatives
 \item \textbf{Medical conditions:} Iron deficiency, Parkinson disease, restless legs syndrome, and occasionally other neurologic or medical illness
\end{itemize}

\section*{Related Symptoms}
\begin{itemize}
 \item Subjective restlessness (inner tension, anxiety, dysphoria)
 \item Objective motor agitation (rocking, shifting, crossing/uncrossing legs, pacing)
 \item Inability to maintain a seated or standing posture
 \item Sleep disturbance and irritability
 \item Dyskinesia, dystonia, or parkinsonism if another drug-induced movement disorder is also present
\end{itemize}
\textbf{Important:} Akathisia is often mistaken for anxiety, agitation, or worsening psychosis. The Barnes Akathisia Rating Scale can support assessment, but diagnosis remains clinical and requires both the person's experience and observed movements when present.

\section*{Medical Evaluation}
\begin{itemize}
 \item Clinical diagnosis based on the subjective urge to move, observed restlessness, functional impact, and the timing of medicine exposure; the Barnes Akathisia Rating Scale may help quantify severity
 \item Medication history with timeline of symptom onset relative to drug initiation, dose change, or withdrawal
 \item Physical examination: observe the person sitting, standing, and walking; check for parkinsonism, dystonia, dyskinesia, tremor, and other neurologic signs
 \item Differentiate from restless legs syndrome (an urge to move mainly in the legs, worse at rest and in the evening, and relieved by movement), anxiety, psychomotor agitation, tardive dyskinesia, and withdrawal
 \item Serum ferritin and iron studies when iron deficiency or restless legs syndrome is suspected; other tests are guided by symptoms rather than required for every case
\end{itemize}

\section*{Treatment Options}
\begin{itemize}
 \item The prescribing clinician may reduce the dose, slow the increase, discontinue the offending medicine, or switch to another treatment such as quetiapine or clozapine when clinically appropriate. Do not make these changes independently because the underlying condition may worsen.
 \item If medication is needed, propranolol or low-dose mirtazapine are commonly considered, after checking for contraindications and interactions. Evidence is limited and the choice should be individualized.
 \item A short course of a benzodiazepine may be considered for severe acute symptoms, but sedation, falls, dependence, and worsening of some intoxications or respiratory disease must be considered.
 \item Anticholinergics such as benztropine are not routine treatment for akathisia; they may be useful when clinically significant parkinsonism or another anticholinergic-responsive movement disorder is also present.
 \item Vitamin B6 has emerging but limited evidence and high doses can cause neuropathy; do not use high-dose pyridoxine without specialist supervision.
 \item Electroconvulsive therapy is an uncommon specialist option for severe, treatment-resistant cases, particularly when a serious underlying psychiatric illness also requires ECT.
\end{itemize}

\section*{Self-Care and Home Management}
\begin{itemize}
 \item Do not stop antipsychotic medication abruptly; consult the prescribing provider first
 \item Avoidance of caffeine, stimulants, and other exacerbating substances
 \item Gentle walking or stretching may help temporarily if it is safe, but do not use strenuous activity to push through severe symptoms
 \item Relaxation techniques: deep breathing, guided imagery, mindfulness
 \item Cool environment and loose clothing to reduce physical discomfort
\end{itemize}

\section*{References}
\begin{thebibliography}{99}
\bibitem{akathisia:ref1} Barnes TRE. ``The Barnes Akathisia Rating Scale -- revisited.'' \textit{J Psychopharmacol}. 2003;17(4):365-370.
\bibitem{akathisia:ref2} Poyurovsky M. ``Acute antipsychotic-induced akathisia revisited.'' \textit{Br J Psychiatry}. 2010;196(2):89-91.
\bibitem{akathisia:ref3} Salem H, Nagpal C, et al. ``Revisiting the management of antipsychotic-induced akathisia.'' \textit{J Clin Psychopharmacol}. 2017;37(3):287-295.
\bibitem{akathisia:ref4} Stahl SM. \textit{Stahl\'s Essential Psychopharmacology}, 5th ed. Cambridge University Press, 2021.
\bibitem{akathisia:ref5} UpToDate. ``Akathisia: clinical features and management.'' Wolters Kluwer, 2023.
\end{thebibliography}
